\subsection{QUOTIENT OF A MANIFOLD BY A LIE GROUP}

Let \(\mathbf{G}\) be a Lie group, \(\mathbf{X}\) a manifold of class \(\mathbf{C}^r\) and \((g,x)\mapsto gx\) a law of left operation (Algebra, Chapter I, § 5, no. 1) of class \(\mathbf{C}^r\) of \(\mathbf{G}\) on \(\mathbf{X}\) . For all \(g\in G\) , let \(\tau (g)\) denote the automorphism \(x\mapsto gx\) of \(\mathbf{X}\) defined by \(g\) . For all \(x\in X\) , let \(\rho (x)\) denote the orbital mapping \(g\mapsto gx\) of \(\mathbf{G}\) into \(\mathbf{X}\) defined by \(x\) . Then

\[
\rho (x) = \rho (g x) \circ \delta (g), \quad \rho (x) = \tau (g) \circ \rho (x) \circ \gamma (g ^ {- 1})\tag{3}
\]

for all \(g\in G\) and \(x\in X\) . Hence

\begin{enumerate}
\def\labelenumi{(\arabic{enumi})}
\setcounter{enumi}{3}
\tightlist
\item
\end{enumerate}

\[
\mathrm{T} _ {g} (\rho (x)) = \mathrm{T} _ {e} (\rho (g x)) \circ \mathrm{T} _ {g} (\delta (g))
\]

\[
\mathrm{T} _ {g} (\rho (x)) = \mathrm{T} _ {x} (\tau (g)) \circ \mathrm{T} _ {e} (\rho (x)) \circ \mathrm{T} _ {g} (\gamma (g ^ {- 1}))\tag{5}
\]

PROPOSITION 9. Let \(x \in X\) and \(g_0 \in G\) .

\begin{enumerate}
\def\labelenumi{(\roman{enumi})}
\item
  If \(\rho(x)\) is an immersion (resp. a submersion, a subimmersion) at \(g_0\) , then, for all \(g \in G\) , \(\rho(gx)\) is an immersion (resp. a submersion, a subimmersion).
\item
  If \(\rho(x)\) is of rank \(k\) at \(g_0\) , then for all \(g \in G\) , \(\rho(gx)\) is of constant rank equal to \(k\) .
\end{enumerate}

This follows immediately from formulae (4) and (5) since \(\mathrm{T}_{g}(\delta(g)),\mathrm{T}_{x}(\tau(g))\) and \(\mathrm{T}_{g}(\gamma(g^{-1}))\) are isomorphisms.

COROLLARY. Let \(x \in \mathbf{X}\) . If \(\mathbf{K}\) is of characteristic 0 and \(\mathbf{X}\) is finite-dimensional, \(\rho(x)\) is a subimmersion. If further \(\rho(x)\) is injective, \(\rho(x)\) is an immersion.

This follows from Proposition 9 and Differentiable and Analytic Manifolds, R, 5.10.6.

Observe that, if \(\eta\) denotes the mapping \((g,x)\mapsto gx\) of \(\mathbf{G}\times \mathbf{X}\) into \(\mathbf{X}\) , then, for \(g\in \mathbf{G},x\in \mathbf{X},u\in \mathrm{T}_g(\mathrm{G}),v\in \mathrm{T}_x(\mathrm{X})\)

\[
\mathrm{T} _ {(g, x)} (\eta) (u, v) = \mathrm{T} _ {(g, x)} (\eta) (u, 0) + \mathrm{T} _ {(g, x)} (\eta) (0, x)
\]

that is

\[
\mathrm{T} _ {(g, x)} (\eta) (u, v) = \mathrm{T} _ {g} (\rho (x)) u + \mathrm{T} _ {x} (\tau (g)) v.\tag{6}
\]

PROPOSITION 10. Let G be a Lie group and X a manifold of class \(\mathbf{C}^r\) with a law of left operation of class \(\mathbf{C}^r\) of G on X. Suppose that:

\begin{enumerate}
\def\labelenumi{(\alph{enumi})}
\item
  the group G operates properly and freely on X;
\item
  for all \(x \in \mathbf{X}\) , \(\rho(x)\) is an immersion (which is a consequence of (a) if \(\mathbf{K}\) is of characteristic 0 and \(\mathbf{X}\) is finite-dimensional).
\end{enumerate}

Then the equivalence relation defined by G on X is regular (Differentiable and Analytic Manifolds, R, 5.9.5). There exists on the quotient set X/G one and only one manifold structure such that the canonical mapping \(\pi: \mathrm{X} \to \mathrm{X} / \mathrm{G}\) is a submersion. The underlying topology of this manifold structure is the quotient topology of that on X; it is Hausdorff. Finally, (X, G, X/G, \(\pi\) ) is a principal left fibre bundle. \(\dagger\)

Let \(\theta\) be the mapping \((g, x) \mapsto (x, gx)\) of \(\mathbf{G} \times \mathbf{X}\) into \(\mathbf{X} \times \mathbf{X}\) . This mapping is of class \(\mathbf{C}^r\) . We show that it is an immersion. For \(u \in \mathrm{T}_g(\mathbf{G})\) and \(v \in \mathrm{T}_x(\mathbf{X})\) , by (6),

\[
\mathrm{T} _ {(g, x)} (\theta) (u, v) = (v, \mathrm{T} _ {g} (\rho (x)) u + \mathrm{T} _ {x} (\tau (g)) v).\tag{7}
\]

But \(\mathrm{T}_{g}(\rho(x))\) is injective by hypothesis (b) and hence \(\mathrm{T}_{(g,x)}(\theta)\) is injective. Its image is the topological direct sum of the subspace \(\mathrm{H}_{g,x}\) consisting of the vectors \((v,\mathrm{T}_x(\tau(g))v)\) for \(v\in \mathrm{T}_x(\mathrm{X})\) and the subspace

\[
\mathrm{I} _ {g, x} = \{0 \} \times \mathrm{T} _ {g} (\rho (x)) (\mathrm{T} _ {g} (\mathrm{G})).
\]

By hypothesis (b), \(\mathrm{T}_g(\rho(x))(\mathrm{T}_g(\mathrm{G}))\) admits a topological supplement \(J_{g,x}\) in \(\mathrm{T}_{gx}(\mathbf{X})\) . Hence the image of \(\mathrm{T}_{(g,x)}(\theta)\) admits the topological supplement \(\{0\} \times J_{g,x}\) . Hence we have proved that \(\theta\) is an immersion of \(\mathbf{G} \times \mathbf{X}\) in \(\mathbf{X} \times \mathbf{X}\) .

As G operates freely on X, \(\theta\) is injective. Let C be the graph of the equivalence relation R defined by G on X. As G operates properly, \(\theta\) is a homeomorphism of \(G \times X\) onto C (General Topology, Chapter I, § 10, Proposition 2). By Differentiable and Analytic Manifolds, R, 5.8.3, it is a submanifold of \(X \times X\) and \(\theta\) is an isomorphism of the manifold \(G \times X\) onto the manifold C. The tangent space \(T_{(x,gx)}(C)\) is identified with

\[
\mathrm{T} _ {(g, x)} (\theta) \left(\mathrm{T} _ {(g, x)} (\mathrm{G} \times \mathrm{X})\right) = \mathrm{H} _ {g, x} \oplus \mathrm{I} _ {g, x} \subset \mathrm{T} _ {(x, g x)} (\mathrm{X} \times \mathrm{X}).
\]

Let \(\mathrm{pr}_1\) and \(\mathrm{pr}_2\) be the canonical projections of \(\mathbf{X} \times \mathbf{X}\) onto the two factors. It is immediate that \(\mathrm{T}_{(x,gx)}(\mathrm{pr}_1)\) maps \(\mathrm{H}_{g,x}\) onto \(\mathrm{T}_x(\mathrm{X})\) and that the kernel of \(\mathrm{T}_{(x,gx)}(\mathrm{pr}_1) \mid \mathrm{T}_{(x,gx)}(\mathrm{C})\) is \(\mathrm{I}_{g,x}\) . Thus \(\mathrm{pr}_1 \mid \mathrm{C}\) is a submersion of \(\mathrm{C}\) onto \(\mathrm{X}\) . By Differentiable and Analytic Manifolds, R, 5.9.5, R is regular. By definition, there therefore exists on the quotient set X/G one and only one manifold structure such that \(\pi\) is a submersion. The underlying topology of X/G is the quotient topology of that of X (Differentiable and Analytic Manifolds, R, 5.9.4). This topology is Hausdorff (General Topology, Chapter III, § 4, Proposition 3).

For all \(b \in \mathrm{X} / \mathrm{G}\) , there exist an open neighbourhood \(\mathrm{W}\) of \(b\) and a morphism \(\sigma: \mathrm{W} \to \mathrm{X}\) such that \(\pi \circ \sigma = \mathrm{id}_{\mathbf{W}}\) (Differentiable and Analytic Manifolds, R, 5.9.1). Let \(\phi\) be the bijection \((g, w) \mapsto g\sigma(w)\) of \(\mathrm{G} \times \mathrm{W}\) onto \(\pi^{-1}(\mathrm{W})\) . It is of class \(C^r\) . Then \(\pi(g\sigma(w)) = w\) and

\[
\theta^ {- 1} (\sigma (w), g \sigma (w)) = (g, \sigma (w))
\]

and hence the inverse bijection of \(\phi\) is of class \(\mathbf{C}^r\) . Clearly \(\phi(gg', w) = g\phi(g', w)\) for \(w \in W\) , \(g \in G\) , \(g' \in G\) . Hence \((\mathbf{X}, \mathbf{G}, \mathbf{X}/\mathbf{G}, \pi)\) is a principal left fibre bundle.

Remark. With the above hypotheses, further let H be a manifold of class \(C^{r}\) and \((x, h) \mapsto m(x, h)\) a mapping of class \(C^{r}\) of X × H into X such that \(m(gx, h) = gm(x, h)\) for \(x \in X, g \in G, h \in H\) . Let n be the mapping of \((\mathrm{X}/\mathrm{G}) \times \mathrm{H}\) into X/G derived from m by taking quotients. We show that n is of class \(C^{r}\) . Consider the diagram

\[
\begin{array}{c} \mathrm{X} \times \mathrm{H} \xrightarrow {m} \mathrm{X} \\ \pi \times 1 \Biggl \downarrow \\ (\mathrm{X} / \mathrm{G}) \times \mathrm{H} \xrightarrow {n} \mathrm{X} / \mathrm{G} \end{array}
\]

It is commutative, \(\pi \circ m\) is of class \(\mathbf{C}^r\) and \(\pi \times 1\) is a surjective submersion; it then suffices to apply Differentiable and Analytic Manifolds, R, 5.9.5.

Let G be a Lie group, X a manifold of class \(C^{r}\) and \((g, x) \mapsto xg\) a law of right operation of class \(C^{r}\) of G on X. Let \(\tau(g)x = \rho(x)g = xg\) for \(g \in G\) , \(x \in X\) . Then this time

\[
\rho (x) = \rho (x g) \circ \gamma (g ^ {- 1}), \quad \rho (x) = \tau (g) \circ \rho (x) \circ \delta (g)\tag{3'}
\]

and hence

(4')

\[
\mathrm{T} _ {g} (\rho (x)) = \mathrm{T} _ {e} (\rho (x g)) \circ \mathrm{T} _ {g} (\gamma (g ^ {- 1}))\tag{5'}
\]

\[
\mathrm{T} _ {g} (\rho (x)) = \mathrm{T} _ {x} (\tau (g)) \circ \mathrm{T} _ {e} (\rho (x)) \circ \mathrm{T} _ {g} (\delta (g)).
\]

On the other hand, if \(\eta\) denotes the mapping \((g, x) \mapsto xg\) of \(G \times X\) into X, formula (6) remains valid. Proposition 9, its Corollary and Proposition 10 remain equally true (with ``principal left fibre bundle'' replaced by ``principal right fibre bundle'' in the last).

